\documentclass[11pt,a4paper]{article}

\usepackage[margin=2.5cm]{geometry}
\usepackage{listings}
\usepackage{xcolor}
\usepackage{hyperref}
\usepackage{titlesec}

\title{mwm: A Minimal Runtime-Configurable X11 Window Manager}
\author{Mohammed Bakr}
\date{\today}

\lstset{
  basicstyle=\ttfamily\small,
  breaklines=true,
  frame=single,
  columns=fullflexible,
  keepspaces=true
}

\begin{document}

\maketitle

\section{Overview}

mwm is a minimal tiling and floating window manager for X11, written in C
using Xlib. It supports both tiling and floating window arrangements,
multiple layout modes, workspaces, and per-application rules. Unlike an
earlier version of this project, configuration is read from a plain text
file at startup rather than compiled into the binary, allowing keybinds,
colors, and behavior to be changed without recompilation.

\section{Design Goals}

\begin{itemize}
  \item Small, readable codebase with no external dependencies beyond Xlib.
  \item Runtime configuration through a simple key-value file format.
  \item Support for both tiling and floating window management, selectable
        per window and per workspace.
  \item Predictable, minimal feature set in the spirit of suckless-style
        window managers.
\end{itemize}

\section{Architecture}

The project is organized into five translation units, each responsible for
one area of functionality.

\begin{itemize}
  \item \textbf{mwm.c}: X server connection, event loop, and window
        lifecycle management (mapping, unmapping, destruction).
  \item \textbf{config.c}: Parses the runtime configuration file and
        populates the global configuration structure, including keybinds
        and window rules.
  \item \textbf{client.c}: Maintains the linked list of managed windows
        (clients), including creation, destruction, lookup, and focus
        handling.
  \item \textbf{layout.c}: Computes window geometry for the tiling and
        monocle layouts. Floating windows retain their own geometry and
        are left untouched by this module.
  \item \textbf{actions.c}: Implements the behavior triggered by each
        keybind, such as spawning programs, changing layouts, and
        switching workspaces.
\end{itemize}

\subsection{Core Data Structures}

The \texttt{Client} structure represents a single managed window. It
stores the window's current geometry, its stored floating geometry (so a
window can return to its previous size and position after being tiled),
its workspace, and its floating and fullscreen state.

The \texttt{Config} structure holds all values loaded from the
configuration file: the modifier key, border appearance, tiling ratio,
gap size, the list of keybinds, and the list of per-application rules.

The \texttt{WM} structure is the single global instance holding the X
display connection, the root window, screen dimensions, the client list,
the currently focused client, and the active layout for each workspace.

\section{Configuration File}

The configuration file is located at \texttt{\textasciitilde/.config/mwm/mwmrc}.
It uses a simple line-based format:

\begin{lstlisting}
key = value
bind = mod+key, action, arg
rule = ClassName, workspace=N, floating=0 or 1
\end{lstlisting}

Lines beginning with \texttt{\#} are treated as comments and ignored by
the parser. If no configuration file is found, built-in defaults are
used instead.

\subsection{General Settings}

\begin{itemize}
  \item \texttt{mod}: the modifier key used for keybinds, either
        \texttt{super} or \texttt{alt}.
  \item \texttt{border\_width}: window border thickness in pixels.
  \item \texttt{border\_focus} and \texttt{border\_unfocus}: border colors
        in hexadecimal, used for the focused and unfocused windows.
  \item \texttt{master\_ratio}: fraction of the screen width given to the
        master area in tiling mode.
  \item \texttt{gap}: spacing in pixels between tiled windows and the
        screen edge.
\end{itemize}

\subsection{Keybinds}

Each \texttt{bind} line associates a key combination with an action and
an optional argument. The key combination is written as a sequence of
modifiers and a key name joined with \texttt{+}, for example
\texttt{mod+shift+Return}. Supported actions include spawning a program,
killing the focused window, moving focus, toggling floating or fullscreen
state, changing the layout, switching or moving windows between
workspaces, resizing the master area, reloading the configuration, and
quitting.

\subsection{Rules}

A \texttt{rule} line assigns a workspace or floating state to windows of
a given class name, applied automatically when the window is first
managed.

\section{Event Handling}

mwm operates on a single-threaded event loop built around
\texttt{XNextEvent}. The events handled are:

\begin{itemize}
  \item \texttt{MapRequest}: a new window is being mapped and is brought
        under management.
  \item \texttt{DestroyNotify} and \texttt{UnmapNotify}: a managed window
        is removed from the client list and the layout is recomputed.
  \item \texttt{KeyPress}: the pressed key combination is matched against
        the configured keybinds, and the corresponding action is invoked.
  \item \texttt{EnterNotify}: used to implement focus-follows-mouse
        behavior.
\end{itemize}

\section{Layout Engine}

Three layout modes are supported, selectable per workspace:

\begin{itemize}
  \item \textbf{Tile}: the first window occupies a master area on the
        left, sized according to \texttt{master\_ratio}. Remaining
        windows are stacked vertically on the right.
  \item \textbf{Monocle}: all windows occupy the full available screen
        area, stacked on top of each other.
  \item \textbf{Float}: no automatic geometry is applied; windows keep
        whatever position and size they were given.
\end{itemize}

Fullscreen windows bypass the active layout and occupy the entire screen
regardless of mode.

\section{Building and Installation}

The project is built with a plain Makefile and requires only the Xlib
development headers.

\begin{lstlisting}
make
sudo make install
\end{lstlisting}

Installation places the binary at \texttt{/usr/local/bin/mwm} and copies
a default configuration file to the user's home directory if one does
not already exist.

\section{Testing}

Because running an unfinished window manager directly in a login session
carries risk, mwm can be tested inside Xephyr, a nested X server that
runs as a window within an existing session.

\begin{lstlisting}
Xephyr -screen 1280x800 :1 &
DISPLAY=:1 ./mwm
\end{lstlisting}

This allows the window manager to be started, interacted with, and
restarted freely without affecting the outer session.

\section{Future Work}

Planned extensions include mouse-based dragging and resizing of floating
windows, an optional status bar, and multi-monitor support.

\end{document}
